\section{Appendix}
\subsection*{Reproduction}
\begin{small}\begin{verbatim}
# the tests (no network, models or services needed)
python -m pytest
# long-conversation benchmark
python -m benchmarks.longconv.run --variants <names> --turns 80 --repeats 3 --keep-chats benchmarks/results/chats
python -m benchmarks.longconv.ladder --chats benchmarks/results/chats --variant full
python -m benchmarks.verifier_audit sample benchmarks/results/ladder_full.json --out audit_sheet.json
# document-QA architecture comparison
python -m benchmarks.docqa.validate
python -m benchmarks.docqa.run --workers 4
python -m benchmarks.docqa.calibrate benchmarks/results/docqa_full75.json --test-from benchmarks/results/docqa_pilot30.json
python -m benchmarks.docqa.report <results.json> --pdf report.pdf --latex tables.tex
\end{verbatim}\end{small}
\subsection*{Hyper-parameters}
\begin{table}[h]\centering\small\begin{tabular}{ll}\toprule Setting & Value \\ \midrule
\texttt{LLM\_MODEL} & gemini-3.5-flash-lite \\
\texttt{LLM\_REASONING\_EFFORT} & low \\
\texttt{EMBEDDING\_MODEL} & BAAI/bge-small-en-v1.5 \\
\texttt{CROSS\_ENCODER\_MODEL} & cross-encoder/ms-marco-MiniLM-L-6-v2 \\
\texttt{PARENT\_CHUNK\_SIZE} & 3000 \\
\texttt{CHILD\_CHUNK\_SIZE} & 800 \\
\texttt{MEMORY\_BUDGET\_TOKENS} & 2000 \\
\texttt{CQM\_W\_GROUNDING} & 0.35 \\
\texttt{CQM\_W\_TOKEN} & 0.25 \\
\texttt{CQM\_W\_COHERENCE} & 0.25 \\
\texttt{CQM\_W\_CONTRADICTION} & 0.15 \\
\texttt{CQM\_TRIGGER\_HEADROOM} & 0.2 \\
\texttt{CQM\_TRIGGER\_CIS} & 0.55 \\
\texttt{STOP\_CONFIDENCE} & 0.8 \\
\texttt{MAX\_EXTRA\_BATCHES} & 3 \\
\texttt{CONF\_W\_MODEL} & 0.5 \\
\texttt{CONF\_W\_SIMILARITY} & 0.25 \\
\texttt{CONF\_W\_SPREAD} & 0.25 \\
\texttt{HISTORY\_TURNS} & 3 \\
\bottomrule\end{tabular}\caption{Settings of the runs (defaults of the code; the values the runs used are in each report's run-settings table).}\end{table}
\subsection*{Probes of \bench{}}
\begin{small}\begin{longtable}{p{0.07\linewidth}p{0.12\linewidth}p{0.75\linewidth}}
\toprule Kind & Bucket & Statement and probe \\ \midrule
recall & near & Note this for later: the draft deadline for my torts coursework is 15 November. \emph{Probe:} What is my draft deadline? \\
recall & mid & For my report I will cover only proximate cause and last clear chance, and I will skip nuisance and defamation. Remember that. \emph{Probe:} Which topics did I say I would skip in my report? \\
recall & far & Keep in mind that my main worked example in the report will be the railway-crossing hypothetical. \emph{Probe:} Which example did I choose as the main one for my report? \\
recall & very_far & Please note that my report has a strict limit of 2,500 words. \emph{Probe:} What word limit did I mention for my report? \\
recall & near & My advisor is Dr. Meera Rao, and she wants Bluebook-style citations. Note that. \emph{Probe:} Who did I say my advisor is, and what citation style does she want? \\
recall & mid & Remember that I am working on this with two classmates, Kabir and Ananya. \emph{Probe:} Which classmates did I say I am working with? \\
recall & far & Note for later: my torts tutorial is every Tuesday at 4 pm. \emph{Probe:} When did I say my torts tutorial is? \\
recall & very_far & Please remember that my viva will be in room B-204. \emph{Probe:} Which room did I say my viva will be in? \\
recall & near & Note this: I keep all my references in Zotero. \emph{Probe:} Which tool did I say I keep my references in? \\
recall & mid & Keep in mind that the report has to be written in British English. \emph{Probe:} Which variety of English did I say the report has to be in? \\
recall & far & Remember that I must also read Prosser's chapter on negligence before the viva. \emph{Probe:} What additional reading did I say I must finish before the viva? \\
recall & very_far & Note for later: I name my laptop Tortoise. \emph{Probe:} What did I say I name my laptop? \\
supersede & mid & Note this: my presentation slot is on 3 December. Later: Correction, please note: my presentation slot has moved to 10 December. \emph{Probe:} When is my presentation slot? \\
supersede & far & Keep in mind that my report topic is the negligence of carriers. Later: Correction, please remember: I changed my report topic to the rescue doctrine. \emph{Probe:} What is my report topic now? \\
absent & -- & (never said) \emph{Probe:} What is my exam date? \\
absent & -- & (never said) \emph{Probe:} What is my student ID number? \\
absent & -- & (never said) \emph{Probe:} Which university did I say I attend? \\
absent & -- & (never said) \emph{Probe:} What grade did I say I got on my last torts assignment? \\
trap & -- & (false premise) \emph{Probe:} What did Romney Marsh v. Trinity House hold about the last clear chance doctrine? \\
trap & -- & (false premise) \emph{Probe:} What did Milwaukee \& St. P. R. Co. v. Kellogg decide about res ipsa loquitur? \\
reference & -- & \emph{Probe (asked at four earlier turns):} Which topic did I ask you to explain at turn {turn}? Answer with the topic only. \\
\bottomrule\end{longtable}\end{small}
\section{Prompts}
\subsection{Snapshot encoder}
\begin{footnotesize}
\begin{verbatim}
You compress part of a conversation about documents into a memory snapshot.
The snapshot replaces these turns later, so anything you leave out is forgotten.

Return ONLY JSON in this shape:
{
  "facts": [{"text": "...", "turn": <turn number>}],
  "user_facts": [{"text": "...", "turn": <turn number>}],
  "entities": [{"name": "...", "turns": [<turn numbers>]}],
  "decisions": [{"topic": "...", "status": "closed" or "open", "turn": <turn number>}],
  "open_questions": ["..."]
}

Rules:
- facts: atomic statements taken from the ANSWERS about the documents. Keep every
    number,
  address, name, date, case name, section number, page and unit exactly as written. Name
    the
  subject in each fact ("Router-2 routes 10.0.0.0/8 via 50.50.50.1"), never "it" or
    "this".
  One statement per fact, no duplicates.
- Skip answers that say information was missing or could not be verified.
- user_facts: things the USER stated in their own messages that are not from a document:
  deadlines, dates, names, limits, preferences, constraints, the example or scope they
    chose,
  corrections to earlier statements ("actually I will also cover X"). Copy the user's
    wording and
  numbers exactly and write each as a full statement ("The user's draft deadline is 15
    November").
- decisions: only choices or commitments the USER made about their own task ("report
    will cover
  only sections 26 and 27, not 32 and 33"); "open" if they are still undecided. Never
    list a topic
  just because it was explained or asked about.
- entities: the named things the facts are about (devices, people, cases, sections,
    concepts).
- open_questions: questions the user asked that the answers did not resolve.
- Use only what the turns say. Do not add outside knowledge.
{exemplars}
TURNS:
{turns}
\end{verbatim}
\end{footnotesize}

\subsection{Reconstruction}
\begin{footnotesize}
\begin{verbatim}
You rebuild conversation memory from a compressed snapshot.

SNAPSHOT of turns {from_turn}-{to_turn} of a conversation about documents:
{snapshot}
{evidence}
TASK: Write what this memory says that is relevant to: {focus}
- Short, self-contained sentences. Name the subject in each; keep numbers, names and
  units exactly as in the snapshot.
- Use the document passages only to restore detail around snapshot facts.
- If nothing is relevant, return an empty list.

Return ONLY JSON: {"sentences": ["...", "..."]}
\end{verbatim}
\end{footnotesize}

\subsection{Sentence judge (verification)}
\begin{footnotesize}
\begin{verbatim}
Check each SENTENCE against the SOURCES below.
For each sentence give the label of the FIRST source (in the order listed) that clearly
states it, or "none" if no source states it or a source contradicts it.
Paraphrases count; a different number, name or claim does not.

{sources}

SENTENCES:
{sentences}

"source" must be exactly one of {labels}: a source label, never a turn or page
    reference.
Return ONLY JSON: {"checks": [{"i": <sentence number>, "source": <label>}]}
\end{verbatim}
\end{footnotesize}

\subsection{Answer verification (confidence check)}
\begin{footnotesize}
\begin{verbatim}
You are an expert RAG verification system.

Your job is to verify whether the generated answer is supported by the retrieved
    context.

Question:
{question}

Retrieved Context:
{context}

Generated Answer:
{answer}

Instructions:

- Compare ONLY with the retrieved context.
- Do NOT use outside knowledge.
- Conclusions that follow directly from the context count as supported, e.g. a path
  traced through connections the context lists, or a result the context shows.
- If the answer is fully supported, return supported=true.
- If any stated fact is missing from or contradicts the context, supported=false.
- A part the answer marks "Not found in the retrieved text" is an honest gap, not an
  unsupported claim. An answer that says information is missing or was never given is
  supported when the context indeed lacks it.
- Check that the answer is about the same thing the question asks about. If the question
  names a device, case, section, person or document and the answer states facts that the
  context attaches to a DIFFERENT one, set "entity_mismatch" to true and supported to
    false.
- An answer that only acknowledges or restates what the user said in the Question
    ("Noted:
  your deadline is 15 November") is supported by the Question itself.
- Sections of the context headed CONVERSATION STATE, HAND-OFF or VERIFIED MEMORY hold
    what
  was said earlier in this conversation; an answer about the conversation may rely on
    them.
- List each unsupported statement in "unsupported_claims" (at most 5, each under 25
  words; empty when supported).
- Return ONLY valid JSON.

Example:

{
    "supported": true,
    "confidence": 0.94,
    "reason": "The answer is fully supported by the retrieved context.",
    "entity_mismatch": false,
    "unsupported_claims": []
}
\end{verbatim}
\end{footnotesize}

\subsection{Baseline prompt: grade (re-implemented method)}
\begin{footnotesize}
\begin{verbatim}
You grade retrieved passages for a question about a novel.

Question: {question}

{passages}

For each passage say "relevant" (it holds information needed to answer), "ambiguous"
    (related, but not clearly
enough) or "irrelevant". Return only JSON: {"grades": ["relevant", "irrelevant"]} with
    one entry per passage, in order.
\end{verbatim}
\end{footnotesize}

\subsection{Baseline prompt: critic (re-implemented method)}
\begin{footnotesize}
\begin{verbatim}
Question about a novel: {question}

Each candidate answer below was written from one passage only.

{candidates}

For every candidate rate from 1 to 5: how relevant the passage is to the question, how
    well the passage supports the answer,
and how useful the answer is. Return only JSON: {"scores": [[relevance, support,
    usefulness], ...]} with one entry per candidate.
\end{verbatim}
\end{footnotesize}

\subsection{Baseline prompt: route (re-implemented method)}
\begin{footnotesize}
\begin{verbatim}
Classify the question about a novel by the retrieval it needs.
A: simple, answerable without looking anything up
B: needs facts from one place in the book (single-step retrieval)
C: needs facts from several places or several reasoning steps (multi-step retrieval)

Question: {question}

Return only the letter.
\end{verbatim}
\end{footnotesize}

\subsection{Baseline prompt: step (re-implemented method)}
\begin{footnotesize}
\begin{verbatim}
Answer the question about a novel from the passages, reasoning step by step.
Write ONLY the next reasoning sentence. If the passages already let you answer fully,
    write ANSWER: and the final answer.

Question: {question}

Passages:
{context}

Reasoning so far: {thoughts}
\end{verbatim}
\end{footnotesize}

\subsection{Baseline prompt: rewrite (re-implemented method)}
\begin{footnotesize}
\begin{verbatim}
Rewrite the question as a short keyword search query for finding the answer in a novel.
    Return only the query.

Question: {question}
\end{verbatim}
\end{footnotesize}
